\documentclass[11pt]{article}

\usepackage[margin=1in]{geometry}
\usepackage{amsmath, amssymb, amsthm}
\usepackage{algorithm}
\usepackage{algpseudocode}
\usepackage{hyperref}
\usepackage{enumitem}

\title{Notes on On-Policy Monte Carlo Control}
\author{}
\date{}

\begin{document}
\maketitle

\section{Setting}

We work in the usual MDP framework: an agent interacts with an environment
through states $s \in \mathcal{S}$, actions $a \in \mathcal{A}$, rewards
$r \in \mathbb{R}$, and a discount factor $\gamma \in [0,1]$. A policy
$\pi(a \mid s)$ is a (possibly stochastic) mapping from states to actions.
The action-value function under $\pi$ is

\[
    Q^\pi(s,a) = \mathbb{E}_\pi\!\left[\, \sum_{k=0}^{\infty} \gamma^k R_{t+k} \;\middle|\; S_t = s, A_t = a \,\right].
\]

Monte Carlo (MC) methods estimate $Q^\pi$ directly from \emph{complete episodes}
of experience, without requiring a model of the environment's dynamics
(unlike dynamic programming). \emph{On-policy} means the policy used to
generate the behaviour (and therefore the data) is the same policy we are
trying to evaluate and improve.

\section{Why not just be greedy?}

If we always act greedily with respect to the current estimate
$\hat{Q}(s,a)$, we may never visit state-action pairs that look bad only
because they have been sampled too few times. MC control needs
\emph{maintained exploration}. The standard fix is an $\varepsilon$-greedy
policy:

\[
    \pi(a \mid s) =
    \begin{cases}
        1 - \varepsilon + \dfrac{\varepsilon}{|\mathcal{A}(s)|}, & a = \arg\max_{a'} \hat{Q}(s,a') \\[4pt]
        \dfrac{\varepsilon}{|\mathcal{A}(s)|}, & \text{otherwise.}
    \end{cases}
\]

This keeps every action selectable with probability at least
$\varepsilon / |\mathcal{A}(s)|$, so every $(s,a)$ pair keeps getting
visited in the limit.

\subsection{GLIE: Greedy in the Limit with Infinite Exploration}

For MC control to provably converge to $Q^*$, the policy sequence
$\{\pi_k\}$ should satisfy the \textbf{GLIE} conditions:

\begin{enumerate}[label=(\roman*)]
    \item Every state-action pair is visited infinitely often.
    \item The policy converges to a greedy policy in the limit,
        e.g.\ $\varepsilon_k \to 0$ as $k \to \infty$ (a common choice is
        $\varepsilon_k = 1/k$).
\end{enumerate}

In practice, many implementations (including the one in this notebook) use
a small \emph{fixed} $\varepsilon$ instead of decaying it. This sacrifices
the formal convergence guarantee but is simpler and works well enough for
small, stochastic-free-ish problems: the learned policy converges to a
near-optimal $\varepsilon$-greedy policy rather than the exactly optimal
greedy one.

\section{Generalized Policy Iteration with MC}

MC control instantiates \textbf{Generalized Policy Iteration (GPI)}: alternate
between

\begin{itemize}
    \item \textbf{Policy evaluation}: estimate $Q^\pi(s,a)$ from sampled
        returns.
    \item \textbf{Policy improvement}: make $\pi$ ($\varepsilon$-)greedy
        with respect to the current $\hat{Q}$.
\end{itemize}

Unlike DP, evaluation here is done from \emph{sampled episodes} rather than
full sweeps over the state space, and the two steps are interleaved at the
granularity of a single episode rather than run to convergence separately.

\section{Computing returns from an episode}

Given an episode
$S_0, A_0, R_1, S_1, A_1, R_2, \dots, S_{T-1}, A_{T-1}, R_T$,
the return from time $t$ is

\[
    G_t = R_{t+1} + \gamma R_{t+2} + \gamma^2 R_{t+3} + \cdots
        = \sum_{k=0}^{T-t-1} \gamma^k R_{t+k+1}.
\]

It is computed efficiently backward through the trajectory:

\[
    G_T = 0, \qquad G_t = R_{t+1} + \gamma\, G_{t+1}.
\]

\section{Every-visit MC control (sample-average updates)}

Let $\text{Returns}(s,a)$ be the (growing) list of returns observed after
visiting $(s,a)$. The value estimate is the sample mean:

\[
    \hat{Q}(s,a) = \frac{1}{N(s,a)} \sum_{i=1}^{N(s,a)} G^{(i)}(s,a).
\]

\begin{algorithm}[h]
\caption{On-policy every-visit MC control ($\varepsilon$-greedy, sample average)}
\begin{algorithmic}[1]
\State Initialize $\hat{Q}(s,a) \leftarrow 0$ for all $s,a$; $\text{Returns}(s,a) \leftarrow \emptyset$
\For{episode $= 1, 2, \dots$}
    \State Generate episode $S_0, A_0, R_1, \dots, S_{T-1}, A_{T-1}, R_T$ using $\varepsilon$-greedy w.r.t.\ $\hat{Q}$
    \State $G \leftarrow 0$
    \For{$t = T-1, T-2, \dots, 0$}
        \State $G \leftarrow R_{t+1} + \gamma G$
        \State append $G$ to $\text{Returns}(S_t, A_t)$
        \State $\hat{Q}(S_t, A_t) \leftarrow \text{mean}\big(\text{Returns}(S_t, A_t)\big)$
    \EndFor
\EndFor
\end{algorithmic}
\end{algorithm}

Note this is \emph{every-visit}, not first-visit: if $(S_t, A_t)$ repeats
within the same episode, each occurrence contributes its own return to the
average. This is exactly what \texttt{on\_policy\_mc\_control} does in
\texttt{Section\_4\_on\_policy\_control\_complete.ipynb}: it stores all
returns per $(s,a)$ pair in a dictionary and recomputes \texttt{np.mean}
after every episode.

A useful incremental identity avoids storing the whole return history: with
$N \leftarrow N+1$,

\[
    \hat{Q}(s,a) \leftarrow \hat{Q}(s,a) + \frac{1}{N(s,a)} \Big( G - \hat{Q}(s,a) \Big).
\]

\section{Constant-$\alpha$ MC control}

Replacing the diminishing step size $1/N(s,a)$ with a fixed
$\alpha \in (0,1]$ gives \textbf{constant-$\alpha$ MC}:

\[
    \hat{Q}(s,a) \leftarrow \hat{Q}(s,a) + \alpha \Big( G - \hat{Q}(s,a) \Big).
\]

\begin{algorithm}[h]
\caption{On-policy MC control, constant step size $\alpha$}
\begin{algorithmic}[1]
\State Initialize $\hat{Q}(s,a) \leftarrow 0$ for all $s,a$
\For{episode $= 1, 2, \dots$}
    \State Generate episode using $\varepsilon$-greedy w.r.t.\ $\hat{Q}$
    \State $G \leftarrow 0$
    \For{$t = T-1, T-2, \dots, 0$}
        \State $G \leftarrow R_{t+1} + \gamma G$
        \State $\hat{Q}(S_t, A_t) \leftarrow \hat{Q}(S_t, A_t) + \alpha \big( G - \hat{Q}(S_t, A_t) \big)$
    \EndFor
\EndFor
\end{algorithmic}
\end{algorithm}

This is the \texttt{constant\_alpha\_mc} function in
\texttt{opmc\_c/Section\_4\_on\_policy\_constant\_alpha\_mc\_complete.ipynb}.

\paragraph{Why use constant $\alpha$?} Two practical reasons:
\begin{itemize}
    \item \textbf{Non-stationarity.} Because the policy keeps changing
        (it is greedy w.r.t.\ a moving $\hat{Q}$), older returns were
        generated under a stale, worse policy. A sample average weighs all
        past returns equally; constant-$\alpha$ instead forms an
        exponentially decaying weighted average, giving recent returns —
        generated under a more recent, better policy — more influence.
    \item \textbf{Memory.} No need to store the full list of past returns
        per $(s,a)$, just the running estimate.
\end{itemize}

The trade-off is that the estimate no longer converges to the true mean
(it keeps tracking / fluctuating), so $\alpha$ acts as a bias-variance /
adaptivity knob: larger $\alpha$ tracks a moving target faster but is
noisier; smaller $\alpha$ is smoother but slower to adapt.

\section{First-visit vs. every-visit}

\begin{itemize}
    \item \textbf{First-visit MC}: only the \emph{first} occurrence of
        $(s,a)$ in an episode contributes a return to the average.
    \item \textbf{Every-visit MC}: \emph{every} occurrence contributes.
\end{itemize}

Both converge to $Q^\pi$ under standard assumptions (each is an unbiased or
asymptotically consistent estimator), but every-visit is simpler to
implement — no need to track "have we seen this pair yet this episode?" —
which is why the notebook implementation above uses it.

\section{Convergence remarks}

\begin{itemize}
    \item MC control with GLIE exploration and sample-average updates
        converges to the optimal $Q^*$ (and hence an optimal
        $\varepsilon$-greedy-turned-greedy policy) with probability 1, as
        both the number of episodes and the exploration schedule satisfy
        the conditions above.
    \item With \emph{fixed} $\varepsilon$ and/or constant $\alpha$, the
        method instead converges (in a suitable stochastic-approximation
        sense) to a stationary distribution around a near-optimal
        $\varepsilon$-greedy policy, trading exact optimality for
        simplicity, adaptivity to non-stationarity, and faster empirical
        learning.
    \item MC methods require \emph{episodic} tasks (no bootstrapping): the
        return $G_t$ can only be computed once the episode terminates.
        This is the key difference from TD-based control methods
        (Sarsa, Q-learning), which bootstrap from $\hat{Q}$ and can learn
        online, step by step, in continuing tasks.
\end{itemize}

\section{Reference}

Sutton, R.\ S., \& Barto, A.\ G.\ \emph{Reinforcement Learning: An
Introduction}, 2nd ed., Chapter 5 (Monte Carlo Methods).

\end{document}
